% Experimental specimen: validate KOMA + CTeX ownership before extracting ipara-core.
% This is NOT a public API example yet.
\documentclass[
  11pt,
  a4paper,
  open=any,
  headinclude=true,
  footinclude=false,
  headheight=18pt,
  footheight=18pt,
  DIV=11
]{scrbook}

% CTeX owns Chinese infrastructure; KOMA owns heading mechanics.
\usepackage[UTF8,fontset=none,heading=false,scheme=chinese]{ctex}

\usepackage{fontspec}
\usepackage{amsmath,mathtools}
\usepackage[warnings-off={mathtools-overbracket,mathtools-colon}]{unicode-math}

\setmainfont{texgyrepagella-regular.otf}[
  BoldFont=texgyrepagella-bold.otf,
  ItalicFont=texgyrepagella-italic.otf,
  BoldItalicFont=texgyrepagella-bolditalic.otf
]
\setsansfont{texgyreheros-regular.otf}[
  Scale=MatchLowercase,
  BoldFont=texgyreheros-bold.otf,
  ItalicFont=texgyreheros-italic.otf,
  BoldItalicFont=texgyreheros-bolditalic.otf
]
\setmonofont{texgyrecursor-regular.otf}[
  Scale=MatchLowercase,
  BoldFont=texgyrecursor-bold.otf,
  ItalicFont=texgyrecursor-italic.otf,
  BoldItalicFont=texgyrecursor-bolditalic.otf
]
\setmathfont{texgyrepagella-math.otf}
\setCJKmainfont{FandolSong-Regular.otf}[
  BoldFont=FandolSong-Bold.otf,
  ItalicFont=FandolKai-Regular.otf
]
\setCJKsansfont{FandolHei-Regular.otf}[
  BoldFont=FandolHei-Bold.otf,
  ItalicFont=FandolKai-Regular.otf
]
\setCJKmonofont{FandolFang-Regular.otf}

\usepackage{xcolor}
\usepackage{booktabs,tabularx,array,longtable}
\usepackage[most]{tcolorbox}
\usepackage{tikz}
\usetikzlibrary{arrows.meta,positioning,calc}
\usepackage{scrlayer-scrpage}
\usepackage{xurl}
\usepackage{microtype}
\usepackage{hyperref}
\usepackage{bookmark}

\definecolor{iparaDeep}{RGB}{42,58,73}
\definecolor{iparaSoft}{RGB}{245,247,249}
\definecolor{iparaRule}{RGB}{145,154,163}

% KOMA owns section/chapter mechanics.
\setkomafont{chapter}{\normalfont\huge\bfseries\sffamily\color{iparaDeep}}
\setkomafont{section}{\normalfont\Large\bfseries\sffamily\color{iparaDeep}}
\setkomafont{subsection}{\normalfont\large\bfseries\sffamily\color{iparaDeep}}
\RedeclareSectionCommand[beforeskip=1.15em,afterskip=0.55em]{section}
\RedeclareSectionCommand[beforeskip=0.9em,afterskip=0.4em]{subsection}
\renewcommand*{\chapterformat}{第\,\thechapter\,章\autodot\enskip}

% KOMA page style only; no fancyhdr.
\clearpairofpagestyles
\ihead{\small\sffamily I.P.A.R.A · Handbook Engine Specimen}
\ohead{\small\sffamily\headmark}
\cfoot{\small\pagemark}
\automark[section]{chapter}
\pagestyle{scrheadings}

\raggedbottom
\setlength{\parindent}{2em}
\setlength{\parskip}{0.22em plus 0.06em minus 0.03em}
\setlength{\emergencystretch}{2em}
\newcolumntype{Y}{>{\raggedright\arraybackslash}X}
\renewcommand{\arraystretch}{1.22}

\newtcolorbox{specimenbox}[1]{
  enhanced,
  breakable,
  colback=iparaSoft,
  colframe=iparaDeep,
  boxrule=0.6pt,
  arc=1.2mm,
  title={#1},
  fonttitle=\bfseries,
  left=1.8mm,right=1.8mm,top=1.2mm,bottom=1.2mm
}

\hypersetup{
  unicode=true,
  colorlinks=true,
  linkcolor=iparaDeep,
  pdftitle={I.P.A.R.A KOMA + CTeX Standard Engine Specimen},
  pdfauthor={I.P.A.R.A}
}

\begin{document}

\frontmatter
\begin{titlepage}
  \centering
  \vspace*{0.18\textheight}
  {\Huge\bfseries\sffamily I.P.A.R.A Handbook Engine\par}
  \vspace{0.7em}
  {\Large Standard Profile · KOMA + CTeX\par}
  \vspace{1.5em}
  {\small Experimental regression specimen — not a public template\par}
\end{titlepage}

\tableofcontents

\mainmatter
\chapter{对象、表示与边界}

\section{为什么先分对象与表示}

设线性映射 $T\colon V\to W$。矩阵 $A$ 只是选定基以后 $T$ 的一种表示；换基改变的是表示，不是映射本身。这个句子同时测试中文、Latin text、数学字体和 KOMA 的 section mark。

English text should remain calm inside a Chinese technical handbook. Inline code such as \texttt{profile=standard} should remain visually distinct without becoming the dominant voice.

\begin{specimenbox}{Mental Model · 表示不是对象}
真正稳定的理解不是记住某个矩阵，而是知道哪些量依赖坐标、哪些结构在合法变换下保持不变。Profile 可以改变这个盒子的 renderer，但 Canonical Source 不应该把知识语义绑定到某种颜色。
\end{specimenbox}

\subsection{公式与长表达式}

短公式：
\[
  \det(P^{-1}AP)=\det A.
\]

较长公式用来观察正文宽度与断行，而不是靠整体缩放挽救：
\begin{align}
  \operatorname{Cov}(X,Y)
  &= \mathbb E\!\left[(X-\mathbb EX)(Y-\mathbb EY)\right] \\
  &= \mathbb E(XY)-\mathbb E(X)\mathbb E(Y).
\end{align}

\section{表格职责}

\begin{center}
\small
\begin{tabularx}{\linewidth}{@{}p{0.19\linewidth}YY@{}}
\toprule
\textbf{层} & \textbf{拥有的东西} & \textbf{不应该拥有的东西}\\
\midrule
KOMA & chapter/section、页眉页脚、TOC 与双面页面 mechanics & Mental Model、Boundary 等学习语义\\
CTeX & 中文语言、CJK 与中文命名基础设施 & Handbook margin 页面状态\\
I.P.A.R.A & Semantic API、Profile 与领域组件 & 重写一套底层 book engine\\
\bottomrule
\end{tabularx}
\end{center}

\section{TikZ 与分页}

\begin{center}
\begin{tikzpicture}[
  node distance=12mm and 14mm,
  box/.style={draw=iparaDeep,rounded corners=1mm,fill=iparaSoft,align=center,minimum height=8mm,inner xsep=6pt},
  edge/.style={-{Latex[length=2mm]},draw=iparaDeep}
]
  \node[box] (a) {KOMA\\mechanics};
  \node[box,right=of a] (b) {CTeX\\Chinese};
  \node[box,below=of $(a)!0.5!(b)$] (c) {I.P.A.R.A\\semantics + profiles};
  \draw[edge] (a) -- (c);
  \draw[edge] (b) -- (c);
\end{tikzpicture}
\end{center}

\chapter{Regression Boundary}

这一章故意产生第二个 chapter，以检查 TOC、chapter mark、页眉更新和页码行为。后续 margin profile 必须复用同一语义内容，而不是复制一份“边栏版正文”。

\section{验收问题}

\begin{enumerate}
  \item 是否出现 CTeX 的 non-standard-class heading warning？
  \item KOMA 是否独立拥有 chapter/section visual hierarchy？
  \item 中文目录名、书签与页眉是否正常？
  \item 是否存在 missing font、undefined reference、unexpected overfull box？
\end{enumerate}

\end{document}
